\documentclass[sigconf,nonacm]{acmart}

\usepackage{booktabs}
\usepackage{listings}
\usepackage{xcolor}
\usepackage{url}

\lstset{
  basicstyle=\ttfamily\small,
  breaklines=true,
  frame=single,
  backgroundcolor=\color{gray!10},
  columns=fullflexible,
}

\begin{document}

\title{FilmSnob: Personalized Movie Recommendations from Letterboxd Reviews}

\author{Rayan Singh}
\affiliation{\institution{University of Illinois Urbana-Champaign}\country{}}
\author{Ryan Vir}
\affiliation{\institution{University of Illinois Urbana-Champaign}\country{}}
\author{Tony Chan}
\affiliation{\institution{University of Illinois Urbana-Champaign}\country{}}

\begin{abstract}
% TODO
\end{abstract}

\maketitle

\section{Introduction}
% TODO

\section{System Overview}

FilmSnob is a content-based movie recommender. We scrape a user's Letterboxd diary (what they watched, how they rated it, what they wrote about it), pull in extra metadata from The Movie Database (TMDB), and build a vector for each film that combines sentence-transformer output with structured features like genre and year. PCA brings the dimensionality down to something manageable, and we recommend films by finding the closest unwatched movies to the user's taste profile via cosine similarity.

The pipeline has four stages:
\begin{enumerate}
  \item \textbf{Letterboxd ingestion:} pull user film diaries via RSS, HTML scraping, or CSV export.
  \item \textbf{TMDB enrichment:} look up plot summaries, genres, cast, director, and keywords through the TMDB API.
  \item \textbf{Cross-user aggregation:} merge per-user data into one film corpus, pooling all reviews for each film.
  \item \textbf{Embedding and recommendation:} encode films with SBERT + metadata, reduce with PCA, and run nearest-neighbor search.
\end{enumerate}

\section{Data Collection}

\subsection{Letterboxd Ingestion}

We support three ways to get data out of Letterboxd, all of which produce the same internal \texttt{UserProfile} object:

\paragraph{RSS feed (primary).}
Every Letterboxd user has a public RSS feed at \texttt{/\{username\}/rss/}. It gives back up to 50 recent diary entries as XML, with custom \texttt{letterboxd:} and \texttt{tmdb:} namespaces. Each entry has the film title, year, the user's star rating (0.5--5.0), watch date, and the TMDB movie ID. Review text shows up in the \texttt{<description>} field as HTML, which we strip out with BeautifulSoup.

We went with RSS over HTML scraping mainly because Letterboxd puts most of its pages behind Cloudflare (we kept hitting HTTP~403s during development). RSS feeds are not blocked. As a bonus, they include the TMDB movie ID directly, so we can skip the title-based search step when looking up metadata.

\paragraph{HTML scraping (fallback).}
We parse the poster grid at each user's films page for slugs and ratings, then their reviews page for review text. This can get a full watch history, but Cloudflare blocks it often enough that it is not reliable as a primary method.

\paragraph{CSV import.}
Letterboxd lets users export their data as CSV files (\texttt{ratings.csv}, \texttt{reviews.csv}, \texttt{watched.csv}). Parsing these gives the most complete history with no scraping involved, but it requires the user to go through the export process themselves.

\subsection{Profile Discovery}

We needed a set of active users to build our film corpus. We scraped the Letterboxd pages for a few popular recent movies (\textit{Oppenheimer}, \textit{Anora}, \textit{The Substance}, etc.), pulled out reviewer usernames from the review links in the page HTML, and checked that each profile page loaded (HTTP~200). This gave us 100 confirmed active users.

\subsection{TMDB Enrichment}

For each film we query the TMDB API v3 detail endpoint with \texttt{append\_to\_response=keywords,credits} and grab:

\begin{itemize}
  \item \textbf{overview:} plot summary (the main text we feed to SBERT)
  \item \textbf{genres:} genre labels
  \item \textbf{cast:} top 10 billed actors
  \item \textbf{director:} pulled from crew credits
  \item \textbf{keywords:} up to 15 thematic tags
  \item \textbf{vote\_average} and \textbf{vote\_count:} TMDB community rating and vote count
\end{itemize}

Most films already have a \texttt{tmdb\_id} from the RSS feed, so we can fetch details by ID directly. For the few that don't, we fall back to searching by title + year, and retry without the year if nothing comes up. We cap API calls at 4 per second to stay under the rate limit.

\subsection{Cross-User Aggregation}

After scraping all 100 users, we merge their data into one corpus. Films are grouped by \texttt{tmdb\_id} so each movie appears once, with all of its user reviews and ratings pooled together. We also store a pre-formatted \texttt{embedding\_text} string (overview + genres + director + keywords) and some aggregate stats like mean rating and review count.

When the same film shows up in multiple users' diaries, we keep all their reviews instead of throwing out duplicates. At embedding time, we encode each review separately with SBERT and average the vectors. A film that 10 people reviewed ends up with a more stable embedding than one with a single review.

\section{Background: PCA}

Principal Component Analysis (PCA) is a standard technique for reducing the number of dimensions in a dataset while keeping as much of the variance as possible. Given a matrix of $n$ data points in $d$ dimensions, PCA finds a new set of axes (the principal components) that are linear combinations of the original features, ordered by how much variance each one explains. The first component points in the direction of greatest variance, the second is orthogonal to the first and captures the next most variance, and so on.

To reduce dimensionality, we keep only the top $k$ components and project every data point onto them, going from $d$ dimensions down to $k$. The key assumption is that if a small number of components explains most of the variance, the remaining dimensions are mostly noise or redundancy and can be dropped without losing much information.

In our case, we start with 791-dimensional film vectors (two 384-d SBERT embeddings plus 23 metadata features). Many of these dimensions are correlated: a review of a horror film will mention similar things as the plot summary, so the content and review embeddings overlap. PCA lets us compress this down to 128 dimensions while retaining the structure that matters for measuring film similarity.

Before running PCA, we standardize each feature to zero mean and unit variance with \texttt{StandardScaler}. This is important because the SBERT dimensions are roughly unit-scale, but features like year and vote count are on completely different scales. Without standardization, PCA would be dominated by whichever features happen to have the largest raw values.

\section{Embedding Construction}

We represent each film as a single vector made up of four groups of features. Before PCA, these add up to 791 dimensions:

\subsection{Content Embedding (384 dimensions)}

We concatenate the plot overview, genre names, director, and keywords into one string (\texttt{embedding\_text}) and run it through the \texttt{all-MiniLM-L6-v2} sentence-transformer (22M parameters, 384-d output). This gives us a 384-dimensional vector for each film's content.

\subsection{Review Embedding (384 dimensions)}

Each user review for a given film is encoded with the same SBERT model, and we average all the review vectors element-wise to get one review embedding per film. Films with no reviews get a zero vector here.

We average in embedding space instead of concatenating the raw review texts. Concatenation would blow past the model's token limit for films with many reviews, and averaging gives each review equal weight no matter how long it is.

\subsection{Scalar Features (4 dimensions)}

\begin{itemize}
  \item \textbf{Year:} release year, so the model can pick up on era preferences
  \item \textbf{Average user rating:} mean Letterboxd star rating across our scraped users
  \item \textbf{TMDB rating:} community vote average (1.0 to 10.0)
  \item \textbf{Log vote count:} $\log(1 + \text{vote\_count})$; we take the log so blockbusters with millions of votes don't dominate
\end{itemize}

\subsection{Genre Multi-Hot (19 dimensions)}

A binary vector with one bit per TMDB genre (19 total). We include this because the SBERT embedding sometimes blurs similar genres together (e.g.\ ``thriller'' and ``horror''), and having explicit genre bits helps keep them distinct.

\subsection{Dimensionality Reduction}

We standardize all 791 features to zero mean and unit variance, then run PCA to bring it down to 128 dimensions. On our corpus, 128 components capture about 62.3\% of the variance. PCA helps mainly because the content and review SBERT vectors overlap a lot for the same film (reviews tend to mention the plot), so there is a lot of redundancy to compress out. It also speeds up the nearest-neighbor search at recommendation time.

\section{Dataset Statistics}

\begin{table}[h]
\begin{tabular}{lr}
\toprule
Metric & Value \\
\midrule
Users scraped & 100 \\
Total film entries (across all users) & 4,852 \\
Unique films (by TMDB ID) & 2,565 \\
Films with TMDB metadata & 2,538 (98.9\%) \\
Films with embedding text & 2,542 (99.1\%) \\
\midrule
User ratings collected & 4,054 (83.6\%) \\
User reviews collected & 3,206 (66.1\%) \\
Films with multiple reviews & 398 (15.5\%) \\
Avg.\ reviews per multi-review film & 4.7 \\
\midrule
Genres represented & 19 \\
Languages represented & 37 \\
Year range & 1896--2026 \\
\midrule
Avg.\ films per user & 48.5 \\
Min / Max films per user & 32 / 50 \\
\bottomrule
\end{tabular}
\caption{Corpus statistics after scraping and enrichment.}
\end{table}

\section{Recommendation Method}

Given a Letterboxd username, the system does the following:

\begin{enumerate}
  \item \textbf{Scrape} the user's RSS feed to get their recent watches, ratings, and reviews.
  \item \textbf{Look up} each film's pre-computed 128-d embedding from the corpus. If a film is not in the corpus, we encode it on the fly with the same pipeline.
  \item \textbf{Build a user profile vector} as a weighted average of their film embeddings:
  \[
    \mathbf{u} = \frac{\sum_{i} r_i \cdot \mathbf{e}_i}{\sum_{i} r_i}
  \]
  where $r_i$ is the star rating for film $i$ and $\mathbf{e}_i$ is its embedding. Films rated higher by the user have more influence on the profile.
  \item \textbf{Rank} every corpus film the user hasn't seen by cosine similarity to $\mathbf{u}$.
  \item \textbf{Return} the top-$N$ results.
\end{enumerate}

There is no training step. Everything runs against pre-computed embeddings, so it is fast. If we add a learned model later, this is the baseline.

\section{Example Recommendations}

We picked three Letterboxd users with pretty different taste to see how the system handles each one. Table~\ref{tab:recs} lists their highest-rated films next to what our system recommended. Because we pull data from each user's RSS feed, which Letterboxd caps at 50 entries, these results are based on their 50 most recent diary entries rather than their full watch history (which may be much larger). Recommendations take under two seconds per user.

\begin{table*}[t]
\small
\begin{tabular}{llll}
\toprule
& \textbf{@davechen} & \textbf{@demyvarda} & \textbf{@broey\_deschanel} \\
& \textit{1,716 films total (50 most recent used)} & \textit{2,861 films total (50 most recent used)} & \textit{287 films total (50 most recent used)} \\
\midrule
\textbf{Top-rated} & Kill Bill: The Whole Bloody Affair (5.0) & Call Me by Your Name (5.0) & Nirvanna the Band the Show the Movie (5.0) \\
\textbf{films}     & Hard Boiled (5.0) & Twin Peaks: Fire Walk with Me (5.0) & If I Had Legs I'd Kick You (5.0) \\
                   & Oldboy (5.0) & Donkey Skin (5.0) & Stories We Tell (5.0) \\
                   & Nirvanna the Band the Show the Movie (5.0) & The Young Girls of Rochefort (5.0) & La Dolce Vita (5.0) \\
\midrule
\textbf{Top}          & Cool Hand Luke (0.50) & The Princess Bride (0.41) & An Unmarried Woman (0.47) \\
\textbf{recs}         & A Clockwork Orange (0.46) & Little Miss Sunshine (0.36) & Sleepless in Seattle (0.46) \\
\textbf{(sim.)}       & Heaven \& Earth (0.45) & Wicked: For Good (0.35) & Something's Gotta Give (0.46) \\
                      & Magnolia (0.44) & Freakier Friday (0.38) & A Summer Place (0.45) \\
                      & Promising Young Woman (0.44) & The Worst Person in the World (0.35) & Late Autumn (0.42) \\
\bottomrule
\end{tabular}
\caption{Recommendations for three users based on their 50 most recent Letterboxd diary entries (the RSS feed limit). Top-rated films and recommendations are ranked by cosine similarity against the user's profile vector.}
\label{tab:recs}
\end{table*}

\paragraph{@davechen} favors violent, stylish genre films (Kill Bill, Hard Boiled, Oldboy). The recs match: Cool Hand Luke, A Clockwork Orange, and Magnolia all share that intense, director-driven quality.

\paragraph{@demyvarda} watches mostly French cinema, musicals, and romance. The recs stay in that lane: Princess Bride, Little Miss Sunshine, The Worst Person in the World. Nothing dark or action-heavy.

\paragraph{@broey\_deschanel} has broad taste (documentary, Italian art cinema, indie drama), but the recs all focus on women's experiences: An Unmarried Woman, Late Autumn, Something's Gotta Give. The system picks up on theme over genre.

All three get completely different lists despite some overlap in what they've watched, since the rating-weighted profile emphasizes what each user likes most.

\section{Reproducibility}

All code is in the project repository. To run the full pipeline from scratch:

\begin{lstlisting}
pip install -r requirements.txt
cp .env.example .env  # add TMDB API key

python scripts/scrape_batch.py
python scripts/build_corpus.py
python scripts/build_embeddings.py
python scripts/recommend.py --username <user>
\end{lstlisting}

\section{Remaining Work}

What we still need to finish before the end of the semester:

\begin{enumerate}
  \item \textbf{Formal evaluation.} Right now we only have the qualitative examples above. We need to split each user's history into train/test, build the profile from training, and see if the test films rank highly. We plan to report Recall@10 and MRR.
  \item \textbf{Expand the corpus beyond 50 films per user.} The RSS feed caps at 50 entries. We want to try the CSV export path or paginated HTML scraping to get full watch histories for at least some users and see if more data improves recommendations.
  \item \textbf{Baseline comparisons.} Compare SBERT + PCA against simpler methods like TF-IDF on plot text, genre-only matching, and popularity ranking. Need to check that the extra complexity actually helps.
  \item \textbf{Rating prediction model.} Train a small MLP or ridge regression that takes a user profile vector and a film embedding and predicts the star rating. This was in the original proposal and would go beyond just nearest-neighbor ranking.
  \item \textbf{Hyperparameter tuning.} Try different PCA dimensions (64, 128, 256), a bigger SBERT model (\texttt{all-mpnet-base-v2}), and different weighting schemes for the user profile vector.
  \item \textbf{User study.} Have a few people try it with their own Letterboxd accounts and get feedback on whether the recommendations are any good.
  \item \textbf{Clean up codebase.} Fill in missing type hints and make sure everything runs end-to-end from a clean clone.
\end{enumerate}

\section{Conclusion}
Summary
The FilmSnob system has both a content-based and predictive model. Given films the user has rated, it retrieves similar films by matching descriptions, genres, casts, keywords using TF-IDF and BM25. It can also predict based on a user’s rating history the score they would assign to any candidate film, and then rank those candidates accordingly.

Data comes from three general sources. The first is movie databases like IMDB and TMDB, whose APIs give us data on any film such as titles, genres, cast, crew, languages, countries of origin, keywords, and more. The second is Kaggle datasets to supply descriptions and keywords to supplement. The third is the Letterboxd platform to provide ratings and reviews from users.

We evaluate the system using personal Letteboxd accounts and mock data. For a given user, we can train the model on half of the ratings and use the model to predict the other half to evaluate it. Comparing predicted ratings to actual ratings provides a direct accuracy measure.

The first few weeks of the project was conducting research on existing models and methods in sentiment analysis, retrieval-augmented generation, and recommendation systems in similar fields. The following weeks were for creating a first draft of the code and creating a report for the first milestone. The final weeks are allocated to model training, integration, testing, and the final written report.

Intro
Letterboxd and MyAnimeList collectively host millions of user-generated film ratings and written reviews. These platforms capture two district signals of user preference which are numerical ratings and natural language reviews that often convey a nuance a score cannot. Despite this, neither platform has leveraged the massive amount of data they have to generate personalized recommendations for users to watch.

FilmSnob is a recommender system that addresses this problem. Given a user’s Letterboxd history, the system performs sentiment analysis on written reviews, weights the resulting sentiment scores against star reviews, and constructs a preference profile for that user. It then retrieves candidate films the user has not seen, ranks them my predicted ratings, and returns a ranked recommendation list.

The central challenge is representing the non-standardized data into something quantifiable. A five-star rating system compresses the user’s thoughts into a number. A written review is better but harder to mine into data. We encode both into a shared embedding vector that gives the model more perspective. We use SBERT for dense semantic embeddings. With this we build a model that estimates the score a user would give unseen films.

Description

Initial Evaluation
For our initial evaluation, we plan on performing a train/test split on individual users’ Letterboxd rating histories in which we will be using around 50% of their ratings for training and the remaining 50% for testing. The predictive model was evaluated by comparing predicted ratings against the actual user rating by using standard regression metrics, something like mean absolute error (MAE) and root mean squared error (RMSE). For our retrieval component, we will assess whether top ranked recommendations aligned with the user’s known preferences in genre and themes. Early results will show that the hybrid approach will perform better than a simple content-based or predictive ranking method, specifically for users with a moderate rating history.

Insights so far
With semantic embedding, we can use extracted data from user reviews to add more meaningful signals when ratings are ambiguous or inconsistent. We also observed that user preferences are often more complex than a simple genre matching in which factors like themes captured in text embeddings are able to play a major role in improving recommendations. Current limitations that we come across include some sarcastic reviews that are inconsistent with users and balancing the weight between numerical ratings and textual ratings which will be a focus for fine tuning the model.


\bibliographystyle{ACM-Reference-Format}

\end{document}
